\subsection{女性性欲的激素周期性波动}

不少女性会发现：一个月里，有时兴致勃勃，有时却完全提不起劲。这未必是"善变"或"不爱了"，而很可能与\textbf{月经周期的激素起伏}有关。本节把这条"激素曲线"讲清楚，也提醒你：\textbf{激素只是影响因素之一，绝不是唯一解释}。

\subsubsection{一个月里的性欲曲线}

\begin{itemize}
	\item \textbf{卵泡期（月经结束到排卵前）}：雌激素逐步升高，身体状态与情绪通常转好，性欲常随之上升。
	\item \textbf{排卵前后}：部分女性的性欲在此期达到高峰。雌激素、促黄体生成素（LH）以及雄激素的短暂波动，被认为与此相关；一种进化假说认为这与生育窗口期的性吸引力增强有关，但证据并不统一。
	\item \textbf{黄体期（排卵后到下次月经前）}：孕激素升高，性欲常下降；经前一周尤为明显——乳房胀痛、腹胀、情绪波动与疲劳都会压低性欲。
	\item \textbf{经期}：因人而异。有人因盆腔充血而性欲升高，有人因不适与顾虑而降低，两种都正常。
\end{itemize}

\subsubsection{背后的激素机制}

\begin{itemize}
	\item \textbf{雌激素}：维持阴道润滑、黏膜健康与性唤起能力；水平过低（如绝经后）会造成干涩与性交痛，间接抑制性欲。
	\item \textbf{孕激素}：多与"性欲下降"相关，尤其在经前与产后阶段。
	\item \textbf{雄激素（睾酮等）}：与\textbf{自发性性欲}（即不需要外部刺激就自然产生的欲望）关系密切。女性的雄激素由卵巢与肾上腺少量分泌，随年龄缓慢下降。
	\item \textbf{其他}：催乳素过高可抑制性腺功能；甲状腺功能异常（亢进或减退）、皮质醇长期升高（慢性压力）都会影响性欲。
\end{itemize}

\subsubsection{口服避孕药与性欲}

\begin{itemize}
	\item 部分女性在服用复方口服避孕药后性欲下降，可能与\textbf{性激素结合球蛋白升高、游离睾酮降低}，或某些孕激素的抗雄激素活性有关。
	\item 也有人没有变化，甚至因不再担心怀孕、经期更规律而性欲回升。\textbf{个体差异极大}。
	\item 若怀疑与避孕药相关，可记录 2--3 个周期后与医生沟通，\textbf{在医生指导下更换配方或剂型观察}，不要自行停药。
\end{itemize}

\subsubsection{围绝经期与绝经后}

\begin{itemize}
	\item 雌激素波动下降，阴道干涩、萎缩与性交痛常见，直接拉低性体验与欲望；局部雌激素或润滑剂可显著改善（详见"润滑剂选择""局部雌激素"）。
	\item 但并非单向下降：部分女性在绝经后因摆脱了怀孕顾虑、关系趋于稳定，性欲与性满意度反而回升。
\end{itemize}

\subsubsection{别把一切都推给激素}

\begin{itemize}
	\item 关系质量、信任与沟通、压力水平、睡眠、情绪（抑郁焦虑）、身体形象与用药，往往\textbf{比激素波动本身影响更大}。
	\item 记录"周期—性欲日记"（同时记录睡眠、压力、关系状态），几周后常能看出自己的规律与真正诱因。
	\item 若持续存在明显困扰，可就医检查性激素、甲状腺功能与催乳素等，再决定是否需要干预。
\end{itemize}

\subsection{女性性欲亢进：被误解的"慕男狂"与真实成因}

"慕男狂"（nymphomania）是一个\textbf{带有严重性别偏见}的旧词，现代医学已弃用。它把女性正常的性需求病理化、污名化。当一位女性性欲明显高于伴侣时，她更可能被指责"不检点"，而非被理解——这种\textbf{双重标准本身，就是一种伤害}。本节讨论真正需要关注的"性欲亢进"，以及它背后的医学与心理成因。

\subsubsection{先分清"高性欲"与"病态亢进"}

\begin{itemize}
	\item \textbf{性欲高是正常差异}：性欲的个体差异极大，有些女性本就欲望较强。只要\textbf{自己可控、不造成痛苦与损害}，就不是病。
	\item \textbf{病态性欲亢进的特征}：性冲动持续过强、\textbf{无法自控}，明知带来负面后果（关系破裂、工作受损、健康风险）仍反复进行，且因此产生明显痛苦或功能损害。
	\item \textbf{关键判断}：不是"欲望有多强"，而是"是否失控、是否带来伤害"。
\end{itemize}

\subsubsection{可能的医学与心理成因}

\begin{itemize}
	\item \textbf{药物}：多巴胺激动剂（治疗帕金森病等）可引起性欲亢进，属已知不良反应；部分抗抑郁药在个别情况下也会引发冲动控制问题。
	\item \textbf{精神疾病}：\textbf{躁狂或轻躁狂发作}（双相情感障碍）常伴性欲亢进、冲动消费与冒险行为；边缘型人格障碍、创伤后应激也与冲动性性行为相关。
	\item \textbf{神经因素}：部分颞叶癫痫、脑部病变可表现为性行为改变。
	\item \textbf{内分泌}：雄激素异常升高、多囊卵巢综合征（PCOS）等可能影响性欲。
	\item \textbf{物质使用}：可卡因、甲基苯丙胺等兴奋剂可诱发强迫性性行为。
	\item \textbf{心理与创伤}：以性行为调节情绪、缓解空虚或强迫性重复创伤模式，也不少见。
\end{itemize}

\subsubsection{诊断定位：不是"性成瘾"}

\begin{itemize}
	\item \textbf{ICD-11} 将其归为\textbf{"强迫性性行为障碍"（Compulsive Sexual Behaviour Disorder, CSBD）}，强调其本质是\textbf{冲动控制障碍}，而非"成瘾"。\textbf{"性成瘾"并非正式医学诊断}，使用时应谨慎。
	\item 评估的核心是：失控、反复、痛苦与功能损害——而不是性行为本身的形式或频率。
\end{itemize}

\subsubsection{如何求助}

\begin{itemize}
	\item \textbf{先排查病因}：内分泌科（激素）、神经内科（排除癫痫/脑病变）、精神心理科（双相、创伤、冲动控制）——很多情况在去除病因后显著改善。
	\item \textbf{心理治疗}：认知行为疗法（CBT）、动机访谈、辩证行为疗法（DBT）可用于改善冲动控制与情绪调节。
	\item \textbf{伴侣与关系支持}：伴侣的理解与共同参与，往往比"道德劝说"有效得多。
	\item \textbf{女性求助的特殊困难}：因羞耻与社会偏见，女性更可能被忽视、被指责或不被相信。请寻找\textbf{专业、非评判}的医生与心理师，不必独自承受。
\end{itemize}

\begin{tcolorbox}[colback=purple!5!white,colframe=purple!55!black,title=贴心小叮咛]
性欲的强弱，本就有很大的正常范围——\textbf{高不是错，低也不是错}。真正需要就医的，是\textbf{失控、痛苦与损害}，而不是"和别人不一样"。请警惕那些用"慕男狂""不检点"给女性贴标签的说法：那不是医学，而是偏见。若你正因此困扰，值得被认真、专业地对待。
\end{tcolorbox}
